
\begin{table}[H]
\centering
\caption{Average Response to Treatments, Median Regression}
\label{tab:a2}
\begin{tabular}{lcc}
\toprule
\textbf{} & \multicolumn{2}{c}{\textbf{Inflation Expectations}}\\ 
\cmidrule(lr){2-3}
& (1) & (2) \\ 
\midrule
\multicolumn{3}{l}{\textbf{Relative to original COPOM statement (Treatment 1)}} \\
T2 & 0.990** & 0.988 \\
& (0.476) & (0.730) \\
\midrule
\multicolumn{3}{l}{\textbf{Relative to original COPOM statement (Treatment 1)}} \\
T3 & -0.025 & 0.275 \\
& (0.643) & (0.768) \\
\midrule
\multicolumn{3}{l}{\textbf{Relative to original G1 article (Treatment 2)}} \\
T4 & 0.430 & 0.602 \\
& (0.461) & (0.615) \\
\midrule
\multicolumn{3}{l}{\textbf{Relative to condensed COPOM statement (Treatment 3)}} \\
T4 & 1.540** & 1.371* \\
& (0.716) & (0.773) \\
\midrule
Demographic Controls & No & Yes \\
Remove Outliers & Yes & Yes \\
\bottomrule
\end{tabular}
\begin{minipage}{\textwidth}
    {\fontsize{8}{8}\selectfont\textit{Notes:} The table reports the average change in inflation expectations of individuals in each treatment group relative to those in the highlighted treatment group. Treatments are described in detail in the text. The second column uses the same specification as the first, but augmented with respondent-specific controls. Results are from Median regressions. Robust standard errors are reported in parenthesis. * p < 0.1; ** p < 0.05; *** p < 0.01.}
\end{minipage}
\end{table}
